% Options for packages loaded elsewhere
\PassOptionsToPackage{unicode}{hyperref}
\PassOptionsToPackage{hyphens}{url}
%
\documentclass[
]{article}
\usepackage{amsmath,amssymb}
\usepackage{iftex}
\ifPDFTeX
  \usepackage[T1]{fontenc}
  \usepackage[utf8]{inputenc}
  \usepackage{textcomp} % provide euro and other symbols
\else % if luatex or xetex
  \usepackage{unicode-math} % this also loads fontspec
  \defaultfontfeatures{Scale=MatchLowercase}
  \defaultfontfeatures[\rmfamily]{Ligatures=TeX,Scale=1}
\fi
\usepackage{lmodern}
\ifPDFTeX\else
  % xetex/luatex font selection
\fi
% Use upquote if available, for straight quotes in verbatim environments
\IfFileExists{upquote.sty}{\usepackage{upquote}}{}
\IfFileExists{microtype.sty}{% use microtype if available
  \usepackage[]{microtype}
  \UseMicrotypeSet[protrusion]{basicmath} % disable protrusion for tt fonts
}{}
\makeatletter
\@ifundefined{KOMAClassName}{% if non-KOMA class
  \IfFileExists{parskip.sty}{%
    \usepackage{parskip}
  }{% else
    \setlength{\parindent}{0pt}
    \setlength{\parskip}{6pt plus 2pt minus 1pt}}
}{% if KOMA class
  \KOMAoptions{parskip=half}}
\makeatother
\usepackage{xcolor}
\usepackage[margin=1in]{geometry}
\usepackage{graphicx}
\makeatletter
\def\maxwidth{\ifdim\Gin@nat@width>\linewidth\linewidth\else\Gin@nat@width\fi}
\def\maxheight{\ifdim\Gin@nat@height>\textheight\textheight\else\Gin@nat@height\fi}
\makeatother
% Scale images if necessary, so that they will not overflow the page
% margins by default, and it is still possible to overwrite the defaults
% using explicit options in \includegraphics[width, height, ...]{}
\setkeys{Gin}{width=\maxwidth,height=\maxheight,keepaspectratio}
% Set default figure placement to htbp
\makeatletter
\def\fps@figure{htbp}
\makeatother
\setlength{\emergencystretch}{3em} % prevent overfull lines
\providecommand{\tightlist}{%
  \setlength{\itemsep}{0pt}\setlength{\parskip}{0pt}}
\setcounter{secnumdepth}{-\maxdimen} % remove section numbering
\usepackage{alex}
\usepackage{xfrac}
\ifLuaTeX
  \usepackage{selnolig}  % disable illegal ligatures
\fi
\IfFileExists{bookmark.sty}{\usepackage{bookmark}}{\usepackage{hyperref}}
\IfFileExists{xurl.sty}{\usepackage{xurl}}{} % add URL line breaks if available
\urlstyle{same}
\hypersetup{
  pdftitle={Time Series Analysis - STAT 478 - Final Exam - Part 1 Q4},
  pdfauthor={Alex Towell (atowell@siue.edu)},
  hidelinks,
  pdfcreator={LaTeX via pandoc}}

\title{Time Series Analysis - STAT 478 - Final Exam - Part 1 Q4}
\author{Alex Towell
(\href{mailto:atowell@siue.edu}{\nolinkurl{atowell@siue.edu}})}
\date{}

\begin{document}
\maketitle

\hypertarget{part-1-problem-4}{%
\section{Part 1: Problem 4}\label{part-1-problem-4}}

Consider the following models where \(\{e_t\}\) is a \(0\) mean white
noise process with variance \(σ^2\).

\begin{enumerate}
\item [i] $Y_t = 1.9 Y_{t−1} − 0.9 Y_{t−2} + e_t − 0.5e_{t−1}$
\item [ii] $Y_t = 0.5 Y_{t−1} + e_t − 0.2e_{t−1} − 0.15e_{t−2}$
\end{enumerate}

\hypertarget{part-a}{%
\subsection{Part (a)}\label{part-a}}

\fbox{Write each model above using backshift notation.}

\begin{enumerate}
\item [(i)] $(1 - 1.9 \backshift + 0.9 \backshift^2) Y_t = (1 - 0.5 \backshift) e_t$.

Working on both sides of the equation simultaneously, we use the following
sequence of transformations to arrive at a canonical form.
\begin{align*}
  (1 - 1.9 \backshift + 0.9 \backshift^2) Y_t &= (1 - \sfrac{1}{2} \backshift) e_t\\
  %%%%%%%%%%%\sfrac{10}{10}(1 - 1.9 \backshift + 0.9 \backshift^2) Y_t &= (1 - \sfrac{1}{2} \backshift) e_t\\
  \sfrac{1}{10}(10 - 19 \backshift + 9 \backshift^2) Y_t &= (1 - \sfrac{1}{2} \backshift) e_t\\
  \sfrac{1}{10}(1-\backshift)(10-9\backshift) Y_t  &=  (1 - \sfrac{1}{2} \backshift) e_t\\
  \sfrac{1}{100}(1-\backshift)(1-\sfrac{9}{10}\backshift) Y_t  &=  (1 - \sfrac{1}{2} \backshift) e_t\\
  (1-\backshift)(1-\sfrac{9}{10}\backshift) Y_t  &=  (1 - \sfrac{1}{2} \backshift) 100 e_t
\end{align*}
We let $\eta_t = 10^2 e_t$, which is white noise with a mean $0$ and variance
$10^4 \sigma^2$, thus
$$
  (1-\backshift)(1-\sfrac{9}{10}\backshift) Y_t = (1 - \sfrac{1}{2} \backshift) \eta_t.
$$

\item [(ii)] $(1 - 0.5 \backshift) Y_t = (1 - 0.2 \backshift - 0.15 \backshift^2) e_t$.
Working on both sides of the equation simultaneously, we use the following
sequence of transformations to arrive at a canonical form.
\begin{align*}
  (1 - 0.5 \backshift) Y_t          &= (1 - 0.2 \backshift - 0.15 \backshift^2) e_t\\
  (1 - \sfrac{1}{2} \backshift) Y_t &= (1 - \sfrac{1}{5} \backshift - \sfrac{3}{20} \backshift^2) e_t\\
  (1 - \sfrac{1}{2} \backshift) Y_t &= (1 - \sfrac{1}{2}\backshift)(1 + \sfrac{3}{10}) e_t\\
                                Y_t &= (1 + \sfrac{3}{10}) e_t
\end{align*}
\end{enumerate}

\hypertarget{part-b}{%
\subsection{Part (b)}\label{part-b}}

\fbox{\begin{minipage}{.9\textwidth}
Characterize these models as models in the $\arima(p,d,q)$ family, that is,
identify $p$, $d$ and $q$.
\end{minipage}}

\begin{enumerate}
\item[(i)] By the form, $\{Y_t\}$ is an $\arima(1,1,1)$ process, or equivalently,
$\{\nabla Y_t\}$ is an $\arma(1,1)$ process.
\item[(ii)] By the form, $\{Y_t\}$ is an $\arima(0,0,1)$ process, or equivalently, $\ma(1)$ process.
\end{enumerate}

\hypertarget{part-c}{%
\subsection{Part (c)}\label{part-c}}

\fbox{Determine if each model corresponds to a stationary process or not.}

\begin{enumerate}
\item[(i)] $\{Y_t\}$ is an $\arima(1,1,1)$ process, which means that its characteristic function has a unit root.
Recall that $\arima(1,1,1)$ denotes a 
$\{Y_t\}$ is therefore a non-stationary process.\footnote{$\nabla \{Y_t\}$ is a stationary $\arma(1,1)$ process since it does not have a unit root.}
\item[(ii)] $\{Y_t\}$ is an $\ma(1)$ process, which are necessarily stationary.
\end{enumerate}

\end{document}
